% Toolbox Vignette of the Finance Toolbox
% Compile with: julia -e 'using PollisTeX; display(compile("finance.tex"))'
% The look of a Pollis Toolbox Vignette (the fonts and colours of the Pollis New Methods documents).
% Set \vignetteauthor, \vignettedate and, if not CC BY-SA 4.0, \vignettelicense after % The look of a Pollis Toolbox Vignette (the fonts and colours of the Pollis New Methods documents).
% Set \vignetteauthor, \vignettedate and, if not CC BY-SA 4.0, \vignettelicense after % The look of a Pollis Toolbox Vignette (the fonts and colours of the Pollis New Methods documents).
% Set \vignetteauthor, \vignettedate and, if not CC BY-SA 4.0, \vignettelicense after \input{vignette_style}.
\documentclass[10pt,a4paper]{article}

\usepackage[a4paper,left=18mm,right=18mm,top=15mm,bottom=17mm,footskip=8mm]{geometry}
\usepackage{iftex}
\ifPDFTeX\usepackage[T1]{fontenc}\fi
\usepackage[default,semibold]{sourcesanspro}
\usepackage{textcomp}
\usepackage[italic]{mathastext}
\usepackage{amsmath}
\usepackage{microtype}
\usepackage[table]{xcolor}
\usepackage{tikz}
\usepackage[most]{tcolorbox}
\usepackage{fontawesome5}
\usepackage{tabularx,booktabs}
\usepackage{enumitem}
\usepackage{fancyhdr}
\usepackage[hidelinks]{hyperref}

% ---------------------------------------------------------------- colours (as in the brochure)
\definecolor{navy}{HTML}{1C3A5E}
\definecolor{pblue}{HTML}{3B78B5}
\definecolor{amber}{HTML}{F2B33D}
\definecolor{deepamber}{HTML}{D9902A}
\definecolor{pale}{HTML}{EEF3F9}
\definecolor{ink}{HTML}{2A2F36}
\definecolor{logonavy}{HTML}{173D63}
\definecolor{logodark}{HTML}{C98A1E}
\definecolor{logoamber}{HTML}{F2B446}
\definecolor{logocream}{HTML}{FEEFD2}
\definecolor{soft}{HTML}{5E6773}
\definecolor{agree}{HTML}{2F7D4F}
\definecolor{differ}{HTML}{B23B3B}
\color{ink}

\setlength{\parindent}{0pt}
\setlength{\parskip}{4pt}
\setlist{nosep,leftmargin=*,itemsep=2pt}
\renewcommand{\arraystretch}{1.15}

\pagestyle{fancy}
\fancyhf{}
\renewcommand{\headrulewidth}{0pt}
\fancyfoot[L]{\footnotesize\color{soft}\textbf{Pollis}\enspace Toolbox Vignette\enspace·\enspace\nmtheme}
\fancyfoot[R]{\footnotesize\color{soft}\thepage}
\newcommand{\nmtheme}{}
\newcommand{\vignetteauthor}{}
\newcommand{\vignettedate}{\today}
\newcommand{\vignettelicense}{\href{https://creativecommons.org/licenses/by-sa/4.0/}{CC BY-SA 4.0}}

% The Pollis logo, redrawn so it prints sharply at any size
\newcommand{\pollislogo}[1]{%
  \begin{tikzpicture}[x=#1/512,y=-#1/512,baseline=(current bounding box.center)]
    \fill[logonavy] (0,0) rectangle (512,512);
    \fill[logodark] (36,422) -- (104,256) -- (172,422) -- cycle;
    \fill[logoamber] (104,422) -- (203,180) -- (302,422) -- cycle;
    \fill[logocream] (203,422) -- (339,89) -- (476,422) -- cycle;
  \end{tikzpicture}}

% Top of the first page: logo, name; then the kicker, title and a one-line promise
\newcommand{\nmhead}[3]{%
  \renewcommand{\nmtheme}{#1}%
  \thispagestyle{fancy}%
  \begin{minipage}[c]{13mm}\pollislogo{11mm}\end{minipage}%
  \begin{minipage}[c]{\dimexpr\linewidth-13mm\relax}
    {\color{navy}\Large\bfseries Pollis}\enspace{\color{soft}Toolbox Vignette}\hfill
    {\color{deepamber}\small\bfseries\MakeUppercase{#1}}
  \end{minipage}\par\vspace{7mm}
  {\color{navy}\fontsize{22}{26}\selectfont\bfseries #2}\par\vspace{4pt}
  {\color{soft}\large #3}\par\vspace{5mm}}

\newcommand{\nmsection}[1]{\par\vspace{5pt}{\color{navy}\large\bfseries #1}\par\vspace{1pt}}
\newcommand{\nmsub}[1]{\par\vspace{2pt}{\color{deepamber}\bfseries #1}\par}
\newcommand{\gt}{\,{\color{deepamber}\bfseries ›}\,}
\newcommand{\menupath}[1]{{\bfseries\color{navy}#1}}
\newcommand{\fn}[1]{\texttt{\small #1}}
% Citations: \nmcite{key}{Breiman (2001)} links to the entry \nmref{key}{...} in the References
\newcommand{\nmcite}[2]{\hyperlink{ref:#1}{\color{pblue}#2}}
\newcommand{\nmref}[2]{\raisebox{1.1em}[0pt][0pt]{\hypertarget{ref:#1}{}}#2\par}
\newcommand{\nmdoi}[1]{\href{https://doi.org/#1}{\color{pblue}doi:#1}}
% Where the package is: its name and a link (its GitHub repository, once published)
\newcommand{\package}[2]{\href{#2}{\fn{#1.jl}}}

% Key facts at the top: source, missing in Julia, package, result
\newcommand{\factsheet}[4]{%
  \begin{tcolorbox}[enhanced, colback=pale, colframe=pale, boxrule=0pt, sharp corners,
    left=7pt, right=7pt, top=5pt, bottom=5pt, fontupper=\small]
  \begin{tabularx}{\linewidth}{@{}>{\leavevmode\bfseries\color{navy}}p{27mm}>{\leavevmode}X@{}}
  Tested against & #1\\
  Missing in Julia & #2\\
  Package & #3\\
  Result & #4
  \end{tabularx}
  \end{tcolorbox}}

% Boxes
\newtcolorbox{pbox}[2][]{enhanced, colback=pale, colframe=amber, boxrule=0pt,
  leftrule=3pt, sharp corners, left=7pt, right=7pt, top=6pt, bottom=6pt,
  fontupper=\small\raggedright, before upper={{\color{navy}\normalsize\bfseries #2}\par\smallskip}, #1}
\newtcolorbox{nbox}[2][]{enhanced, colback=navy, colframe=navy, boxrule=0pt,
  sharp corners, left=8pt, right=8pt, top=7pt, bottom=7pt, coltext=white,
  fontupper=\small\raggedright, before upper={{\color{amber}\normalsize\bfseries #2}\par\smallskip}, #1}

% Value tables: published beside Julia
\newenvironment{valuetable}[1]{%
  \par\small\tabularx{\linewidth}{#1}\toprule}{\bottomrule\endtabularx\par}
\newcommand{\ok}{{\color{agree}\faCheck}}
\newcommand{\no}{{\color{differ}\faTimes}}

% End of the document: credit and licence
\newcommand{\nmcredit}{%
  \par\vfill
  \begin{minipage}[b]{0.78\linewidth}\footnotesize\color{soft}
  A Pollis Toolbox Vignette, written from the panels of the toolbox.\par
  \textcopyright\ \vignetteauthor\enspace·\enspace Licensed under \vignettelicense\enspace·\enspace\vignettedate
  \end{minipage}\hfill
  \begin{minipage}[b]{0.18\linewidth}\raggedleft\pollislogo{9mm}\end{minipage}}
.
\documentclass[10pt,a4paper]{article}

\usepackage[a4paper,left=18mm,right=18mm,top=15mm,bottom=17mm,footskip=8mm]{geometry}
\usepackage{iftex}
\ifPDFTeX\usepackage[T1]{fontenc}\fi
\usepackage[default,semibold]{sourcesanspro}
\usepackage{textcomp}
\usepackage[italic]{mathastext}
\usepackage{amsmath}
\usepackage{microtype}
\usepackage[table]{xcolor}
\usepackage{tikz}
\usepackage[most]{tcolorbox}
\usepackage{fontawesome5}
\usepackage{tabularx,booktabs}
\usepackage{enumitem}
\usepackage{fancyhdr}
\usepackage[hidelinks]{hyperref}

% ---------------------------------------------------------------- colours (as in the brochure)
\definecolor{navy}{HTML}{1C3A5E}
\definecolor{pblue}{HTML}{3B78B5}
\definecolor{amber}{HTML}{F2B33D}
\definecolor{deepamber}{HTML}{D9902A}
\definecolor{pale}{HTML}{EEF3F9}
\definecolor{ink}{HTML}{2A2F36}
\definecolor{logonavy}{HTML}{173D63}
\definecolor{logodark}{HTML}{C98A1E}
\definecolor{logoamber}{HTML}{F2B446}
\definecolor{logocream}{HTML}{FEEFD2}
\definecolor{soft}{HTML}{5E6773}
\definecolor{agree}{HTML}{2F7D4F}
\definecolor{differ}{HTML}{B23B3B}
\color{ink}

\setlength{\parindent}{0pt}
\setlength{\parskip}{4pt}
\setlist{nosep,leftmargin=*,itemsep=2pt}
\renewcommand{\arraystretch}{1.15}

\pagestyle{fancy}
\fancyhf{}
\renewcommand{\headrulewidth}{0pt}
\fancyfoot[L]{\footnotesize\color{soft}\textbf{Pollis}\enspace Toolbox Vignette\enspace·\enspace\nmtheme}
\fancyfoot[R]{\footnotesize\color{soft}\thepage}
\newcommand{\nmtheme}{}
\newcommand{\vignetteauthor}{}
\newcommand{\vignettedate}{\today}
\newcommand{\vignettelicense}{\href{https://creativecommons.org/licenses/by-sa/4.0/}{CC BY-SA 4.0}}

% The Pollis logo, redrawn so it prints sharply at any size
\newcommand{\pollislogo}[1]{%
  \begin{tikzpicture}[x=#1/512,y=-#1/512,baseline=(current bounding box.center)]
    \fill[logonavy] (0,0) rectangle (512,512);
    \fill[logodark] (36,422) -- (104,256) -- (172,422) -- cycle;
    \fill[logoamber] (104,422) -- (203,180) -- (302,422) -- cycle;
    \fill[logocream] (203,422) -- (339,89) -- (476,422) -- cycle;
  \end{tikzpicture}}

% Top of the first page: logo, name; then the kicker, title and a one-line promise
\newcommand{\nmhead}[3]{%
  \renewcommand{\nmtheme}{#1}%
  \thispagestyle{fancy}%
  \begin{minipage}[c]{13mm}\pollislogo{11mm}\end{minipage}%
  \begin{minipage}[c]{\dimexpr\linewidth-13mm\relax}
    {\color{navy}\Large\bfseries Pollis}\enspace{\color{soft}Toolbox Vignette}\hfill
    {\color{deepamber}\small\bfseries\MakeUppercase{#1}}
  \end{minipage}\par\vspace{7mm}
  {\color{navy}\fontsize{22}{26}\selectfont\bfseries #2}\par\vspace{4pt}
  {\color{soft}\large #3}\par\vspace{5mm}}

\newcommand{\nmsection}[1]{\par\vspace{5pt}{\color{navy}\large\bfseries #1}\par\vspace{1pt}}
\newcommand{\nmsub}[1]{\par\vspace{2pt}{\color{deepamber}\bfseries #1}\par}
\newcommand{\gt}{\,{\color{deepamber}\bfseries ›}\,}
\newcommand{\menupath}[1]{{\bfseries\color{navy}#1}}
\newcommand{\fn}[1]{\texttt{\small #1}}
% Citations: \nmcite{key}{Breiman (2001)} links to the entry \nmref{key}{...} in the References
\newcommand{\nmcite}[2]{\hyperlink{ref:#1}{\color{pblue}#2}}
\newcommand{\nmref}[2]{\raisebox{1.1em}[0pt][0pt]{\hypertarget{ref:#1}{}}#2\par}
\newcommand{\nmdoi}[1]{\href{https://doi.org/#1}{\color{pblue}doi:#1}}
% Where the package is: its name and a link (its GitHub repository, once published)
\newcommand{\package}[2]{\href{#2}{\fn{#1.jl}}}

% Key facts at the top: source, missing in Julia, package, result
\newcommand{\factsheet}[4]{%
  \begin{tcolorbox}[enhanced, colback=pale, colframe=pale, boxrule=0pt, sharp corners,
    left=7pt, right=7pt, top=5pt, bottom=5pt, fontupper=\small]
  \begin{tabularx}{\linewidth}{@{}>{\leavevmode\bfseries\color{navy}}p{27mm}>{\leavevmode}X@{}}
  Tested against & #1\\
  Missing in Julia & #2\\
  Package & #3\\
  Result & #4
  \end{tabularx}
  \end{tcolorbox}}

% Boxes
\newtcolorbox{pbox}[2][]{enhanced, colback=pale, colframe=amber, boxrule=0pt,
  leftrule=3pt, sharp corners, left=7pt, right=7pt, top=6pt, bottom=6pt,
  fontupper=\small\raggedright, before upper={{\color{navy}\normalsize\bfseries #2}\par\smallskip}, #1}
\newtcolorbox{nbox}[2][]{enhanced, colback=navy, colframe=navy, boxrule=0pt,
  sharp corners, left=8pt, right=8pt, top=7pt, bottom=7pt, coltext=white,
  fontupper=\small\raggedright, before upper={{\color{amber}\normalsize\bfseries #2}\par\smallskip}, #1}

% Value tables: published beside Julia
\newenvironment{valuetable}[1]{%
  \par\small\tabularx{\linewidth}{#1}\toprule}{\bottomrule\endtabularx\par}
\newcommand{\ok}{{\color{agree}\faCheck}}
\newcommand{\no}{{\color{differ}\faTimes}}

% End of the document: credit and licence
\newcommand{\nmcredit}{%
  \par\vfill
  \begin{minipage}[b]{0.78\linewidth}\footnotesize\color{soft}
  A Pollis Toolbox Vignette, written from the panels of the toolbox.\par
  \textcopyright\ \vignetteauthor\enspace·\enspace Licensed under \vignettelicense\enspace·\enspace\vignettedate
  \end{minipage}\hfill
  \begin{minipage}[b]{0.18\linewidth}\raggedleft\pollislogo{9mm}\end{minipage}}
.
\documentclass[10pt,a4paper]{article}

\usepackage[a4paper,left=18mm,right=18mm,top=15mm,bottom=17mm,footskip=8mm]{geometry}
\usepackage{iftex}
\ifPDFTeX\usepackage[T1]{fontenc}\fi
\usepackage[default,semibold]{sourcesanspro}
\usepackage{textcomp}
\usepackage[italic]{mathastext}
\usepackage{amsmath}
\usepackage{microtype}
\usepackage[table]{xcolor}
\usepackage{tikz}
\usepackage[most]{tcolorbox}
\usepackage{fontawesome5}
\usepackage{tabularx,booktabs}
\usepackage{enumitem}
\usepackage{fancyhdr}
\usepackage[hidelinks]{hyperref}

% ---------------------------------------------------------------- colours (as in the brochure)
\definecolor{navy}{HTML}{1C3A5E}
\definecolor{pblue}{HTML}{3B78B5}
\definecolor{amber}{HTML}{F2B33D}
\definecolor{deepamber}{HTML}{D9902A}
\definecolor{pale}{HTML}{EEF3F9}
\definecolor{ink}{HTML}{2A2F36}
\definecolor{logonavy}{HTML}{173D63}
\definecolor{logodark}{HTML}{C98A1E}
\definecolor{logoamber}{HTML}{F2B446}
\definecolor{logocream}{HTML}{FEEFD2}
\definecolor{soft}{HTML}{5E6773}
\definecolor{agree}{HTML}{2F7D4F}
\definecolor{differ}{HTML}{B23B3B}
\color{ink}

\setlength{\parindent}{0pt}
\setlength{\parskip}{4pt}
\setlist{nosep,leftmargin=*,itemsep=2pt}
\renewcommand{\arraystretch}{1.15}

\pagestyle{fancy}
\fancyhf{}
\renewcommand{\headrulewidth}{0pt}
\fancyfoot[L]{\footnotesize\color{soft}\textbf{Pollis}\enspace Toolbox Vignette\enspace·\enspace\nmtheme}
\fancyfoot[R]{\footnotesize\color{soft}\thepage}
\newcommand{\nmtheme}{}
\newcommand{\vignetteauthor}{}
\newcommand{\vignettedate}{\today}
\newcommand{\vignettelicense}{\href{https://creativecommons.org/licenses/by-sa/4.0/}{CC BY-SA 4.0}}

% The Pollis logo, redrawn so it prints sharply at any size
\newcommand{\pollislogo}[1]{%
  \begin{tikzpicture}[x=#1/512,y=-#1/512,baseline=(current bounding box.center)]
    \fill[logonavy] (0,0) rectangle (512,512);
    \fill[logodark] (36,422) -- (104,256) -- (172,422) -- cycle;
    \fill[logoamber] (104,422) -- (203,180) -- (302,422) -- cycle;
    \fill[logocream] (203,422) -- (339,89) -- (476,422) -- cycle;
  \end{tikzpicture}}

% Top of the first page: logo, name; then the kicker, title and a one-line promise
\newcommand{\nmhead}[3]{%
  \renewcommand{\nmtheme}{#1}%
  \thispagestyle{fancy}%
  \begin{minipage}[c]{13mm}\pollislogo{11mm}\end{minipage}%
  \begin{minipage}[c]{\dimexpr\linewidth-13mm\relax}
    {\color{navy}\Large\bfseries Pollis}\enspace{\color{soft}Toolbox Vignette}\hfill
    {\color{deepamber}\small\bfseries\MakeUppercase{#1}}
  \end{minipage}\par\vspace{7mm}
  {\color{navy}\fontsize{22}{26}\selectfont\bfseries #2}\par\vspace{4pt}
  {\color{soft}\large #3}\par\vspace{5mm}}

\newcommand{\nmsection}[1]{\par\vspace{5pt}{\color{navy}\large\bfseries #1}\par\vspace{1pt}}
\newcommand{\nmsub}[1]{\par\vspace{2pt}{\color{deepamber}\bfseries #1}\par}
\newcommand{\gt}{\,{\color{deepamber}\bfseries ›}\,}
\newcommand{\menupath}[1]{{\bfseries\color{navy}#1}}
\newcommand{\fn}[1]{\texttt{\small #1}}
% Citations: \nmcite{key}{Breiman (2001)} links to the entry \nmref{key}{...} in the References
\newcommand{\nmcite}[2]{\hyperlink{ref:#1}{\color{pblue}#2}}
\newcommand{\nmref}[2]{\raisebox{1.1em}[0pt][0pt]{\hypertarget{ref:#1}{}}#2\par}
\newcommand{\nmdoi}[1]{\href{https://doi.org/#1}{\color{pblue}doi:#1}}
% Where the package is: its name and a link (its GitHub repository, once published)
\newcommand{\package}[2]{\href{#2}{\fn{#1.jl}}}

% Key facts at the top: source, missing in Julia, package, result
\newcommand{\factsheet}[4]{%
  \begin{tcolorbox}[enhanced, colback=pale, colframe=pale, boxrule=0pt, sharp corners,
    left=7pt, right=7pt, top=5pt, bottom=5pt, fontupper=\small]
  \begin{tabularx}{\linewidth}{@{}>{\leavevmode\bfseries\color{navy}}p{27mm}>{\leavevmode}X@{}}
  Tested against & #1\\
  Missing in Julia & #2\\
  Package & #3\\
  Result & #4
  \end{tabularx}
  \end{tcolorbox}}

% Boxes
\newtcolorbox{pbox}[2][]{enhanced, colback=pale, colframe=amber, boxrule=0pt,
  leftrule=3pt, sharp corners, left=7pt, right=7pt, top=6pt, bottom=6pt,
  fontupper=\small\raggedright, before upper={{\color{navy}\normalsize\bfseries #2}\par\smallskip}, #1}
\newtcolorbox{nbox}[2][]{enhanced, colback=navy, colframe=navy, boxrule=0pt,
  sharp corners, left=8pt, right=8pt, top=7pt, bottom=7pt, coltext=white,
  fontupper=\small\raggedright, before upper={{\color{amber}\normalsize\bfseries #2}\par\smallskip}, #1}

% Value tables: published beside Julia
\newenvironment{valuetable}[1]{%
  \par\small\tabularx{\linewidth}{#1}\toprule}{\bottomrule\endtabularx\par}
\newcommand{\ok}{{\color{agree}\faCheck}}
\newcommand{\no}{{\color{differ}\faTimes}}

% End of the document: credit and licence
\newcommand{\nmcredit}{%
  \par\vfill
  \begin{minipage}[b]{0.78\linewidth}\footnotesize\color{soft}
  A Pollis Toolbox Vignette, written from the panels of the toolbox.\par
  \textcopyright\ \vignetteauthor\enspace·\enspace Licensed under \vignettelicense\enspace·\enspace\vignettedate
  \end{minipage}\hfill
  \begin{minipage}[b]{0.18\linewidth}\raggedleft\pollislogo{9mm}\end{minipage}}

\renewcommand{\vignetteauthor}{Pollis}
\renewcommand{\vignettedate}{11 November 2026}
\renewcommand{\vignettelicense}{AGPL-3.0}

\begin{document}

\nmhead{Finance Toolbox}
  {Finance Toolbox}
  {What its panels are for, and when to use each.}

Installing the toolbox adds the submenu \menupath{Extend \gt Finance}, with these panels. Each has an example you can edit and run, and actions on its results.

\medskip
{\small
\begin{tabularx}{\linewidth}{@{}>{\raggedright\arraybackslash}p{62mm}X@{}}
\toprule
\textbf{Panel} & \textbf{Use it when} \\\midrule
Measuring Volatility & You have daily open, high, low and close prices and want volatility from them. \\
High-Frequency Volatility & You have the trades or quotes of one asset within a day and want its volatility, with or without its jumps. \\
High-Frequency Covariance & You have intraday prices of several assets that trade at different times and want their covariance or correlation. \\
ARIMAX Models & One series is forecast from its own past, with or without a seasonal pattern or other variables that drive it. \\
GARCH-Type Models & Volatility clusters and you want it day by day, its forecast, Value-at-Risk or changing correlations. \\
Regime Switching Models & A relationship seems to change between calm and turbulent periods. \\
Point Processes & Events arrive at random times, such as trades, orders or defaults, and you want their rate, whether they excite one another, or simulated ones. \\
Advanced Time-Series Forecasting & A series has several seasonal periods or many zeros, or you need a hard-to-beat benchmark. \\
Neural Network Time-Series Forecasting & Several series mix short patterns with long cycles, and a simple forecast is not enough. \\
Automatic Time-Series Forecasting & You want a forecast without choosing the model by hand. \\
Bayesian Linear Stochastic Discount Factor & There are many candidate factors and the question is which of them price assets. \\
Fama-MacBeth Regression & You want the risk premia of a few factors from a set of test assets. \\
Portfolio Optimisation & You want weights for a set of assets and to know whether they beat equal weights on later returns. \\
Loss Risk Measures & You want the drawdowns, downside risk, Value at Risk or Expected Shortfall of a return series. \\
Trading Agents & You want to see how traders' beliefs and the order book itself produce spreads, fat tails and price impact. \\\bottomrule
\end{tabularx}}


\nmsection{Measuring Volatility}
Seven estimators from \package{RangeVol}{https://github.com/Trumpingtons/RangeVol.jl}. The panel's title is Range-Based Volatility Estimation. The example simulates 500 days of one-minute prices, with a fifth of the variance overnight; replace them with your own. The efficiencies below are relative to close-to-close, under the conditions each estimator assumes.

\medskip
{\small
\begin{tabularx}{\linewidth}{@{}>{\raggedright\arraybackslash}p{36mm}>{\raggedright\arraybackslash}p{58mm}X@{}}
\toprule
\textbf{Estimator} & \textbf{Use it for} & \textbf{What it gives} \\\midrule
Close-to-Close & Only closing prices, or a baseline & The standard deviation of log-returns \\
Parkinson & 24-hour markets with little drift & The high-low range of \nmcite{parkinson1980}{Parkinson (1980)}, about 5 times as efficient \\
Garman-Klass & Full prices, little drift, no overnight gaps & The range and the open-to-close return of \nmcite{garman1980}{Garman and Klass (1980)}, about 7 times \\
Rogers-Satchell & Trending prices without large gaps & The drift-independent estimator of \nmcite{rogers1991}{Rogers and Satchell (1991)} \\
GK-Yang-Zhang & An exchange that closes overnight & Garman-Klass plus the overnight jump \\
Yang-Zhang & Both drift and gaps, the usual case for equities & The minimum-variance blend of \nmcite{yang2000}{Yang and Zhang (2000)}, about 14 times \\
Alizadeh-Brandt-Diebold & A state-space model of log-volatility & The log range, nearly Gaussian around log-volatility (\nmcite{alizadeh2002}{Alizadeh, Brandt and Diebold 2002}) \\\bottomrule
\end{tabularx}}

\medskip
Empirical Efficiency measures the gain on your own data, and Drift and Overnight Gaps checks the assumptions that separate the estimators. Each estimator gives one number for its window, so Volatility Clustering shows when a single full-sample figure is misleading.

\nmsection{High-Frequency Analysis}
Two panels in \menupath{Extend \gt Finance \gt High-Frequency Analysis}, both on \package{HighFrequencyCovariance}{https://github.com/s-baumann/HighFrequencyCovariance.jl} (\nmcite{baumann2022}{Baumann and Klymak 2022}), High-Frequency Volatility also on \package{TwoScalesVolatility}{https://github.com/Trumpingtons/TwoScalesVolatility.jl}. Their examples simulate the ticks; replace them with a table of time, asset and log price.

High-Frequency Volatility measures one asset over a day of ticks. Realized sums squared returns at a sampling interval (\nmcite{andersen1998}{Andersen and Bollerslev 1998}), Two-Scales uses every tick and corrects for microstructure noise, with the estimator, its small-sample adjustment and its rule for the number of sub-grids of \nmcite{zhang2005}{Zhang, Mykland and A\"{\i}t-Sahalia (2005)}, and Bipower leaves jumps out (\nmcite{barndorff2004}{Barndorff-Nielsen and Shephard 2004}). The actions draw the volatility signature plot, estimate the variance of the noise by the same two-scales method, test for jumps after \nmcite{barndorff2006}{Barndorff-Nielsen and Shephard (2006)} and give the jump share of the variance.

High-Frequency Covariance estimates the covariance of several assets that trade at different times: realized covariance on a time grid, pre-averaging (\nmcite{christensen2013}{Christensen, Podolskij and Vetter 2013}), two scales (\nmcite{zhang2011}{Zhang 2011}) and the realized kernel (\nmcite{barndorff2011}{Barndorff-Nielsen, Hansen, Lunde and Shephard 2011}). The actions check whether the matrix is positive semi-definite, repair it when it is not, and give the volatility of a portfolio.

\nmsection{ARIMAX Models}
Three tabs. The panel's title is ARIMA Models. \package{StateSpaceModels}{https://github.com/LAMPSPUC/StateSpaceModels.jl} fits ARIMA, where differencing removes a trend or a drifting level, and Seasonal ARIMA, for data with a repeating period. ARIMAX, with \package{Durbyn}{https://github.com/taf-society/Durbyn.jl}, regresses the series on other variables with ARIMA errors. Forecasting with it needs the future values of those variables, and its intervals leave out any uncertainty about them. Its backtest uses the values that actually occurred, so it scores better than a real forecast would when they have to be predicted too. The actions check the residuals, run an ADF unit root test, and backtest. For the first two tabs they also select the orders by AICc and simulate 500 future paths; for ARIMAX, With and Without the Regressors compares the fit.

\nmsection{GARCH-Type Models}
\package{ARCHModels}{https://github.com/s-broda/ARCHModels.jl} estimates the conditional volatility of returns, following \nmcite{engle1982}{Engle (1982)}. GARCH is the symmetric form of \nmcite{bollerslev1986}{Bollerslev (1986)}. EGARCH (\nmcite{nelson1991}{Nelson 1991}) and GJR-GARCH (\nmcite{glosten1993}{Glosten, Jagannathan and Runkle 1993}) let bad news raise volatility more than good news of the same size, which the News Impact Curve shows. DCC (\nmcite{engle2002}{Engle 2002}) adds changing correlations across several return series. The actions compare innovation distributions and lag orders, backtest Value-at-Risk, and report persistence and half-life.

\nmsection{Regime Switching Models}
\package{MarSwitching}{https://github.com/m-dadej/MarSwitching.jl} lets coefficients and variance change with an unobserved regime that follows a Markov chain, as in \nmcite{hamilton1989}{Hamilton (1989)}. Markov-Switching Regression gives one set of coefficients per regime, and Markov-Switching Autoregression one set of autoregressive dynamics. Markov-Switching ARCH lets volatility cluster within each regime. Time-Varying Transitions makes the switching probabilities depend on an indicator (\nmcite{filardo1994}{Filardo 1994}). Regimes are inferred, not observed, so check the estimates with Stability Across Starting Values and the number of regimes with AIC and BIC.

\nmsection{Point Processes}
The submenu \menupath{Point Processes} has two panels, Model and Simulate, for events that arrive at random times: trades and orders in high-frequency data, or defaults that spread from one firm to another (\nmcite{bacry2015}{Bacry, Mastromatteo and Muzy 2015}). Both cover five processes. Poisson has a constant rate, Inhomogeneous a rate that changes with time, and Cox a rate that is itself random. Hawkes is self-exciting: each event raises the rate for a while, as in \nmcite{hawkes1971}{Hawkes (1971)}, and its branching ratio, below 1 for a stable process, is the share of events triggered by earlier ones. Network Hawkes lets the events of one series excite the others along the links of a network. Model fits a process to event times with \package{PointProcesses}{https://github.com/JoseKling/PointProcesses.jl}, \package{HawkesProcesses}{https://github.com/dm13450/HawkesProcesses.jl} and \package{NetworkHawkesProcesses}{https://github.com/cswaney/NetworkHawkesProcesses.jl}; the Time-Rescaling Test and its Q-Q Plot check the fit (\nmcite{ogata1988}{Ogata 1988}), and Next Event Time and Expected Count predict. Simulate draws events from a process with the parameters given, and Count vs Theory and Multiple Seeds compare them with what the process should give.

\nmsection{Advanced Time-Series Forecasting}
Three methods from \package{Durbyn}{https://github.com/taf-society/Durbyn.jl}. BATS and TBATS (\nmcite{delivera2011}{De Livera, Hyndman and Snyder 2011}) handle several or non-integer seasonal periods, such as hourly data with daily and weekly cycles. The Theta Method (\nmcite{assimakopoulos2000}{Assimakopoulos and Nikolopoulos 2000}) is a simple benchmark for monthly or yearly series. Intermittent Demand is for sparse series with many zeros: \nmcite{croston1972}{Croston (1972)} smooths demand sizes and intervals separately, and SBA and SBJ correct its upward bias, SBA being the correction of \nmcite{syntetos2005}{Syntetos and Boylan (2005)}. Demand Pattern classifies a series before you choose, and Holdout Accuracy scores the forecast on data the model has not seen.

\nmsection{Neural Network Time-Series Forecasting}
Four architectures from \package{FluxArchitectures}{https://github.com/sdobber/FluxArchitectures.jl}, on \package{Flux}{https://github.com/FluxML/Flux.jl}. The package pins Flux 0.13, so the panel runs on that older version. The example forecasts one exchange rate six steps ahead from several, training on the first windows and testing on later ones.

\medskip
{\small
\begin{tabularx}{\linewidth}{@{}>{\raggedright\arraybackslash}p{22mm}>{\raggedright\arraybackslash}p{58mm}X@{}}
\toprule
\textbf{Model} & \textbf{Structure} & \textbf{Use it for} \\\midrule
LSTnet & CNN, GRU, skip connections and AR (\nmcite{lai2018}{Lai, Chang, Yang and Liu 2018}) & Short local patterns mixed with long seasonal cycles \\
DA-RNN & Dual-stage attention (\nmcite{qin2017}{Qin, Song, Chen, Cheng, Jiang and Cottrell 2017}) & Attention over both the driving series and the time steps \\
TPA-LSTM & LSTM with temporal pattern attention (\nmcite{shih2019}{Shih, Sun and Lee 2019}) & Many variables with different periodicities \\
DSANet & Global and local self-attention (\nmcite{huang2019}{Huang, Wang, Wu and Tang 2019}) & Many correlated series, in parallel and with no recurrence \\\bottomrule
\end{tabularx}}

\medskip
Run the Naive Baseline action first: a network that does not beat it adds nothing.

\nmsection{Automatic Time-Series Forecasting}
\package{Durbyn}{https://github.com/taf-society/Durbyn.jl} chooses the model for you. Auto ARIMA sets the differencing by unit-root tests and the orders by a stepwise search, as in \nmcite{hyndman2008}{Hyndman and Khandakar (2008)}. Auto ETS chooses the form of exponential smoothing by information criterion (\nmcite{hyndman2002}{Hyndman, Koehler, Snyder and Grose 2002}). Auto ARARMA applies the memory-shortening filter of \nmcite{parzen1982}{Parzen (1982)} and then a short ARMA, and suits slowly decaying series with no clear seasonal period. Model Race fits them all and keeps the one most accurate on the last observations. Selection Stability shows whether the chosen model changes as the sample does, and Rolling-Origin Accuracy gives a fairer score than a single holdout.

\nmsection{Bayesian Linear Stochastic Discount Factor}
\package{BayesianFactorZoo}{https://github.com/eohne/BayesianFactorZoo.jl}, the method of \nmcite{bryzgalova2023}{Bryzgalova, Huang and Julliard (2023)}. Bayesian SDF gives the full posterior of risk prices, in Sharpe-ratio units. Spike-and-Slab adds the probability that each factor is included, to select factors, and Spike-and-Slab (Traded) also prices tradable factors as test assets. By default the example simulates 600 months of five factors, of which two earn a premium, and 25 test assets; give it your own matrices to replace them. Results come from MCMC, so read Trace Plots and Effective Sample Size before the Risk Price Table.

\nmsection{Fama-MacBeth Regression}
Risk premia from test assets, in two passes as in \nmcite{fama1973}{Fama and MacBeth (1973)}, using \package{BayesianFactorZoo}{https://github.com/eohne/BayesianFactorZoo.jl}. Classical estimates betas over time and then OLS and GLS premia across assets, with the errors-in-variables correction of \nmcite{shanken1992}{Shanken (1992)}; it suits a few established factors with strong betas. Bayesian re-runs the cross-sectional regression on every posterior draw, so a useless factor gets a wide posterior and not a falsely tight estimate, the problem shown by \nmcite{kan1999}{Kan and Zhang (1999)}. The example simulates a market factor, a value factor and a useless one, and Weak Factor Check shows the difference.

\nmsection{Portfolio Optimisation}
Weights for a set of assets from their daily returns, long-only and fully invested, using \package{PortfolioOptimisers}{https://github.com/dcelisgarza/PortfolioOptimisers.jl} with the \package{Clarabel}{https://github.com/oxfordcontrol/Clarabel.jl} solver. Mean-Variance trades expected return against variance, after \nmcite{markowitz1952}{Markowitz (1952)}, and Mean-Expected Shortfall against the average loss on the worst days, minimised by linear programming as in \nmcite{rockafellar2000}{Rockafellar and Uryasev (2000)}; both offer the minimum-risk, the maximum risk-adjusted return and the maximum utility portfolio. HRP, the Hierarchical Risk Parity of \nmcite{lopezdeprado2016}{Lopez de Prado (2016)}, and HERC, the Hierarchical Equal Risk Contribution of \nmcite{raffinot2018}{Raffinot (2018)}, allocate by risk along a tree of the assets built from their correlations, with no expected returns. The actions refit each method on a rolling two-year window and hold it for a quarter, and compare the Sharpe ratios on the held-out returns with those of 1/N, the benchmark that \nmcite{demiguel2009}{DeMiguel, Garlappi and Uppal (2009)} found hard to beat.

\nmsection{Loss Risk Measures}
Measures of the losses of one return series, in four tabs, using \package{Distributions}{https://github.com/JuliaStats/Distributions.jl} and \package{StatsBase}{https://github.com/JuliaStats/StatsBase.jl}. Drawdown gives the maximum and average drawdown with the Calmar and Sterling ratios, and Downside Risk the semi-deviation, the downside deviation below a minimum acceptable return, and the Sortino and Omega ratios, after \nmcite{sortino1994}{Sortino and Price (1994)}. Value at Risk is estimated by historical simulation, from the normal distribution and with the expansion of \nmcite{cornish1938}{Cornish and Fisher (1938)}, and Expected Shortfall, a coherent risk measure in the sense of \nmcite{artzner1999}{Artzner, Delbaen, Eber and Heath (1999)} and \nmcite{acerbi2002}{Acerbi and Tasche (2002)}, by the same three methods. VaR Backtest counts the breaches of rolling out-of-sample forecasts and tests their rate with an exact binomial test, the small-sample form of the test of \nmcite{kupiec1995}{Kupiec (1995)}; \nmcite{christoffersen1998}{Christoffersen (1998)} also tests whether breaches cluster. The estimates are unconditional or from a moving window, so they react slowly when volatility changes.

\nmsection{Trading Agents}
Three simulated markets, each matched through a limit order book. Trading Agents follows \nmcite{chiarella2002}{Chiarella and Iori (2002)} on \package{LimitOrderBook}{https://github.com/p-casgrain/LimitOrderBook.jl}: each period one agent mixes a fundamentalist, a chartist and a noise forecast and sends a one-share order, with mean weights of 10, 1.2 and 1 as in \nmcite{chiarella2009}{Chiarella, Iori and Perell\'{o} (2009)}. Limit Order Book is the zero-intelligence order flow of \nmcite{smith2003}{Smith, Farmer, Gillemot and Krishnamurthy (2003)}, which shows what the book alone does to the spread and to price impact. Market Microstructure runs a live \package{Brokerage}{https://github.com/aaron-wheeler/Brokerage.jl} exchange with \package{TradingAgents}{https://github.com/aaron-wheeler/TradingAgents.jl} traders and a market maker for a session of a few minutes in real time. The actions compare scenarios over 20 seeds, with stronger chartists, without chartists or without fundamentalists, and set the order book against a market that clears at one price. Its agents trade one share with no budget, so a market led by chartists can run away from the fundamental value.

\nmsection{References}
{\footnotesize\sloppy
\nmref{acerbi2002}{Acerbi, C. and Tasche, D. (2002). Expected shortfall: a natural coherent alternative to value at risk. \emph{Economic Notes} 31(2), 379--388. \nmdoi{10.1111/1468-0300.00091}}
\nmref{alizadeh2002}{Alizadeh, S., Brandt, M. W. and Diebold, F. X. (2002). Range-based estimation of stochastic volatility models. \emph{The Journal of Finance} 57(3), 1047--1091. \nmdoi{10.1111/1540-6261.00454}}
\nmref{andersen1998}{Andersen, T. G. and Bollerslev, T. (1998). Answering the skeptics: yes, standard volatility models do provide accurate forecasts. \emph{International Economic Review} 39(4), 885--905. \nmdoi{10.2307/2527343}}
\nmref{artzner1999}{Artzner, P., Delbaen, F., Eber, J.-M. and Heath, D. (1999). Coherent measures of risk. \emph{Mathematical Finance} 9(3), 203--228. \nmdoi{10.1111/1467-9965.00068}}
\nmref{assimakopoulos2000}{Assimakopoulos, V. and Nikolopoulos, K. (2000). The theta model: a decomposition approach to forecasting. \emph{International Journal of Forecasting} 16(4), 521--530. \nmdoi{10.1016/S0169-2070(00)00066-2}}
\nmref{bacry2015}{Bacry, E., Mastromatteo, I. and Muzy, J.-F. (2015). Hawkes processes in finance. \emph{Market Microstructure and Liquidity} 1, 1550005. \nmdoi{10.1142/S2382626615500057}}
\nmref{barndorff2004}{Barndorff-Nielsen, O. E. and Shephard, N. (2004). Power and bipower variation with stochastic volatility and jumps. \emph{Journal of Financial Econometrics} 2(1), 1--37. \nmdoi{10.1093/jjfinec/nbh001}}
\nmref{barndorff2006}{Barndorff-Nielsen, O. E. and Shephard, N. (2006). Econometrics of testing for jumps in financial economics using bipower variation. \emph{Journal of Financial Econometrics} 4(1), 1--30. \nmdoi{10.1093/jjfinec/nbi022}}
\nmref{barndorff2011}{Barndorff-Nielsen, O. E., Hansen, P. R., Lunde, A. and Shephard, N. (2011). Multivariate realised kernels: consistent positive semi-definite estimators of the covariation of equity prices with noise and non-synchronous trading. \emph{Journal of Econometrics} 162(2), 149--169. \nmdoi{10.1016/j.jeconom.2010.07.009}}
\nmref{baumann2022}{Baumann, S. and Klymak, M. (2022). HighFrequencyCovariance: a Julia package for estimating covariance matrices using high frequency financial data. \emph{Journal of Statistical Software} 103(14). \nmdoi{10.18637/jss.v103.i14}}
\nmref{bollerslev1986}{Bollerslev, T. (1986). Generalized autoregressive conditional heteroskedasticity. \emph{Journal of Econometrics} 31(3), 307--327. \nmdoi{10.1016/0304-4076(86)90063-1}}
\nmref{bryzgalova2023}{Bryzgalova, S., Huang, J. and Julliard, C. (2023). Bayesian solutions for the factor zoo: We just ran two quadrillion models. \emph{The Journal of Finance} 78(1), 487--557. \nmdoi{10.1111/jofi.13197}}
\nmref{chiarella2002}{Chiarella, C. and Iori, G. (2002). A simulation analysis of the microstructure of double auction markets. \emph{Quantitative Finance} 2(5), 346--353. \nmdoi{10.1088/1469-7688/2/5/303}}
\nmref{chiarella2009}{Chiarella, C., Iori, G. and Perell\'{o}, J. (2009). The impact of heterogeneous trading rules on the limit order book and order flows. \emph{Journal of Economic Dynamics and Control} 33(3), 525--537. \nmdoi{10.1016/j.jedc.2008.08.001}}
\nmref{christensen2013}{Christensen, K., Podolskij, M. and Vetter, M. (2013). On covariation estimation for multivariate continuous It\^{o} semimartingales with noise in non-synchronous observation schemes. \emph{Journal of Multivariate Analysis} 120, 59--84. \nmdoi{10.1016/j.jmva.2013.05.002}}
\nmref{christoffersen1998}{Christoffersen, P. F. (1998). Evaluating interval forecasts. \emph{International Economic Review} 39(4), 841--862. \nmdoi{10.2307/2527341}}
\nmref{cornish1938}{Cornish, E. A. and Fisher, R. A. (1938). Moments and cumulants in the specification of distributions. \emph{Revue de l'Institut International de Statistique} 5(4), 307--320. \nmdoi{10.2307/1400905}}
\nmref{croston1972}{Croston, J. D. (1972). Forecasting and stock control for intermittent demands. \emph{Journal of the Operational Research Society} 23(3), 289--303. \nmdoi{10.1057/jors.1972.50}}
\nmref{delivera2011}{De Livera, A. M., Hyndman, R. J. and Snyder, R. D. (2011). Forecasting time series with complex seasonal patterns using exponential smoothing. \emph{Journal of the American Statistical Association} 106(496), 1513--1527. \nmdoi{10.1198/jasa.2011.tm09771}}
\nmref{demiguel2009}{DeMiguel, V., Garlappi, L. and Uppal, R. (2009). Optimal versus naive diversification: how inefficient is the 1/N portfolio strategy? \emph{Review of Financial Studies} 22(5), 1915--1953. \nmdoi{10.1093/rfs/hhm075}}
\nmref{engle1982}{Engle, R. F. (1982). Autoregressive conditional heteroscedasticity with estimates of the variance of United Kingdom inflation. \emph{Econometrica} 50(4), 987--1007. \nmdoi{10.2307/1912773}}
\nmref{engle2002}{Engle, R. (2002). Dynamic conditional correlation: A simple class of multivariate generalized autoregressive conditional heteroskedasticity models. \emph{Journal of Business \& Economic Statistics} 20(3), 339--350. \nmdoi{10.1198/073500102288618487}}
\nmref{fama1973}{Fama, E. F. and MacBeth, J. D. (1973). Risk, return, and equilibrium: Empirical tests. \emph{Journal of Political Economy} 81(3), 607--636. \nmdoi{10.1086/260061}}
\nmref{filardo1994}{Filardo, A. J. (1994). Business-cycle phases and their transitional dynamics. \emph{Journal of Business \& Economic Statistics} 12(3), 299--308. \nmdoi{10.1080/07350015.1994.10524545}}
\nmref{garman1980}{Garman, M. B. and Klass, M. J. (1980). On the estimation of security price volatilities from historical data. \emph{The Journal of Business} 53(1), 67--78. \nmdoi{10.1086/296072}}
\nmref{glosten1993}{Glosten, L. R., Jagannathan, R. and Runkle, D. E. (1993). On the relation between the expected value and the volatility of the nominal excess return on stocks. \emph{The Journal of Finance} 48(5), 1779--1801. \nmdoi{10.1111/j.1540-6261.1993.tb05128.x}}
\nmref{hamilton1989}{Hamilton, J. D. (1989). A new approach to the economic analysis of nonstationary time series and the business cycle. \emph{Econometrica} 57(2), 357--384. \nmdoi{10.2307/1912559}}
\nmref{hawkes1971}{Hawkes, A. G. (1971). Spectra of some self-exciting and mutually exciting point processes. \emph{Biometrika} 58(1), 83--90. \nmdoi{10.1093/biomet/58.1.83}}
\nmref{huang2019}{Huang, S., Wang, D., Wu, X. and Tang, A. (2019). DSANet: Dual self-attention network for multivariate time series forecasting. In \emph{Proceedings of the 28th ACM International Conference on Information and Knowledge Management}, 2129--2132. \nmdoi{10.1145/3357384.3358132}}
\nmref{hyndman2008}{Hyndman, R. J. and Khandakar, Y. (2008). Automatic time series forecasting: The forecast package for R. \emph{Journal of Statistical Software} 27(3). \nmdoi{10.18637/jss.v027.i03}}
\nmref{hyndman2002}{Hyndman, R. J., Koehler, A. B., Snyder, R. D. and Grose, S. (2002). A state space framework for automatic forecasting using exponential smoothing methods. \emph{International Journal of Forecasting} 18(3), 439--454. \nmdoi{10.1016/S0169-2070(01)00110-8}}
\nmref{kan1999}{Kan, R. and Zhang, C. (1999). Two-pass tests of asset pricing models with useless factors. \emph{The Journal of Finance} 54(1), 203--235. \nmdoi{10.1111/0022-1082.00102}}
\nmref{kupiec1995}{Kupiec, P. H. (1995). Techniques for verifying the accuracy of risk measurement models. \emph{The Journal of Derivatives} 3(2), 73--84. \nmdoi{10.3905/jod.1995.407942}}
\nmref{lai2018}{Lai, G., Chang, W.-C., Yang, Y. and Liu, H. (2018). Modeling long- and short-term temporal patterns with deep neural networks. In \emph{The 41st International ACM SIGIR Conference on Research \& Development in Information Retrieval}, 95--104. \nmdoi{10.1145/3209978.3210006}}
\nmref{lopezdeprado2016}{Lopez de Prado, M. (2016). Building diversified portfolios that outperform out of sample. \emph{Journal of Portfolio Management} 42(4), 59--69. \nmdoi{10.3905/jpm.2016.42.4.059}}
\nmref{markowitz1952}{Markowitz, H. M. (1952). Portfolio selection. \emph{Journal of Finance} 7(1), 77--91. \nmdoi{10.2307/2975974}}
\nmref{nelson1991}{Nelson, D. B. (1991). Conditional heteroskedasticity in asset returns: A new approach. \emph{Econometrica} 59(2), 347--370. \nmdoi{10.2307/2938260}}
\nmref{ogata1988}{Ogata, Y. (1988). Statistical models for earthquake occurrences and residual analysis for point processes. \emph{Journal of the American Statistical Association} 83(401), 9--27. \nmdoi{10.1080/01621459.1988.10478560}}
\nmref{parkinson1980}{Parkinson, M. (1980). The extreme value method for estimating the variance of the rate of return. \emph{The Journal of Business} 53(1), 61--65. \nmdoi{10.1086/296071}}
\nmref{parzen1982}{Parzen, E. (1982). ARARMA models for time series analysis and forecasting. \emph{Journal of Forecasting} 1(1), 67--82. \nmdoi{10.1002/for.3980010108}}
\nmref{qin2017}{Qin, Y., Song, D., Chen, H., Cheng, W., Jiang, G. and Cottrell, G. W. (2017). A dual-stage attention-based recurrent neural network for time series prediction. In \emph{Proceedings of the Twenty-Sixth International Joint Conference on Artificial Intelligence}, 2627--2633. \nmdoi{10.24963/ijcai.2017/366}}
\nmref{raffinot2018}{Raffinot, T. (2018). The hierarchical equal risk contribution portfolio. SSRN working paper. \nmdoi{10.2139/ssrn.3237540}}
\nmref{rockafellar2000}{Rockafellar, R. T. and Uryasev, S. (2000). Optimization of conditional value-at-risk. \emph{Journal of Risk} 2(3), 21--41. \nmdoi{10.21314/JOR.2000.038}}
\nmref{rogers1991}{Rogers, L. C. G. and Satchell, S. E. (1991). Estimating variance from high, low and closing prices. \emph{The Annals of Applied Probability} 1(4), 504--512. \nmdoi{10.1214/aoap/1177005835}}
\nmref{shanken1992}{Shanken, J. (1992). On the estimation of beta-pricing models. \emph{Review of Financial Studies} 5(1), 1--33. \nmdoi{10.1093/rfs/5.1.1}}
\nmref{shih2019}{Shih, S.-Y., Sun, F.-K. and Lee, H.-y. (2019). Temporal pattern attention for multivariate time series forecasting. \emph{Machine Learning} 108(8--9), 1421--1441. \nmdoi{10.1007/s10994-019-05815-0}}
\nmref{smith2003}{Smith, E., Farmer, J. D., Gillemot, L. and Krishnamurthy, S. (2003). Statistical theory of the continuous double auction. \emph{Quantitative Finance} 3(6), 481--514. \nmdoi{10.1088/1469-7688/3/6/307}}
\nmref{sortino1994}{Sortino, F. A. and Price, L. N. (1994). Performance measurement in a downside risk framework. \emph{The Journal of Investing} 3(3), 59--64. \nmdoi{10.3905/joi.3.3.59}}
\nmref{syntetos2005}{Syntetos, A. A. and Boylan, J. E. (2005). The accuracy of intermittent demand estimates. \emph{International Journal of Forecasting} 21(2), 303--314. \nmdoi{10.1016/j.ijforecast.2004.10.001}}
\nmref{yang2000}{Yang, D. and Zhang, Q. (2000). Drift-independent volatility estimation based on high, low, open, and close prices. \emph{The Journal of Business} 73(3), 477--492. \nmdoi{10.1086/209650}}
\nmref{zhang2005}{Zhang, L., Mykland, P. A. and A\"{\i}t-Sahalia, Y. (2005). A tale of two time scales: determining integrated volatility with noisy high-frequency data. \emph{Journal of the American Statistical Association} 100(472), 1394--1411. \nmdoi{10.1198/016214505000000169}}
\nmref{zhang2011}{Zhang, L. (2011). Estimating covariation: Epps effect, microstructure noise. \emph{Journal of Econometrics} 160(1), 33--47. \nmdoi{10.1016/j.jeconom.2010.03.012}}
}

\nmcredit
\end{document}
